\section{Conclusion}
\label{sec:conclusion}

This study changed the question for river water-quality telemetry over a one-way satellite link from how to deliver the whole record to what the server can soundly decide from what arrives. A certificate language turns each piece of received information into an interval and answers threshold, extremum, conductance-change, and rolling-mean queries with exceeded, not exceeded, or undetermined; the answers are sound when the intervals are valid, exact only without witnesses, and limited by the tolerance for readings that lie near a threshold.

The method was confirmed twice on sealed cohorts. The first confirmation failed its integrity hypothesis and used a comparator that discarded the record; the revised method was confirmed on 35 sealed monitoring sites at byte level, where all three hypotheses that could fail held. No certain decision was wrong. At the unconfigured thresholds of the study, targeted witnesses decided 0.18 to 0.23 more hard cases than a node sending the configured answers with its record, and 0.20 to 0.28 more than the record alone, at a cost of 0.03 to 0.08 in whole-record recovery. The design was adopted although that cost exceeded its own written adoption limit at six messages per day, and the paper reports it as such. Nodes that send the configured answers remained better on the thresholds they encode.

The claim of this study is soundness, and more decisions on hard cases at thresholds the node was not configured for, chiefly thresholds more extreme than the configured ones; it is not the best decisions for known thresholds, and the evidence is trace-driven simulation. Two steps would move it toward practice: a measurement of real channel loss and error statistics, and an implementation of the frame on a satellite IoT module. The software, the frozen plans, and every result file are in the project repository.
